% ---
% title: Three Terms in Cambridge
% date: 2026-07
% category:
% nodes: reflections, work
% links:
% ---
\documentclass{article}
\usepackage[utf8]{inputenc}
\usepackage{amsmath, amssymb, amsthm}
\usepackage{geometry}
\usepackage{graphicx}
\usepackage{hyperref}

\geometry{a4paper, margin=1in}

\title{Three Terms in Cambridge}
\date{July 2026}
\author{Daksh Mehta}

\begin{document}

This week marks the unofficial end of Easter term at Cambridge, exams are finished and I hear the undergraduates cheering outside. As the year comes to an end and all the medicine I've forgotten comes crashing back into my day to day life, I'd like to reflect on some of the things I learnt and did. I'm not finished yet but it feels like a natural point to sit back and think. Most of these reflections will seem obvious and reading back, I find myself questioning why I was surprised by any of this, but I guess no matter how much you read about methods or statistical tests there is a certain lived experience with making the decisions for yourself. Before I get into it, I would like to say how amazing the scientific community at Cambridge is, I attended so many lectures, events and courses and met some extraordinary researchers who took the time to talk to me about their work. I am also very grateful to Professor Bays, the \href{https://bayslab.com}{Computational Cognition Lab} and the Department of Psychology for facilitating my studies this year. 

There was something uniquely challenging about this year, especially for me. I experienced a sort of discomfort sitting with the same problem for so long. I am more used to effort under uncertainty that follows with release. At different time scales: question followed by an answer, writing and feedback, research questions with quick answers found in retrospective data, designing a tool and seeing how well it works in the environment it was made for. It feels like I have been stuck in a year long question that has missing words within it and no options. Whilst this may not be the most novel revelation – everyone knows research is an endeavour to answer questions that no one has answered yet. I felt that year-long focus required a type of stamina I hadn't trained before. 

The work was mostly fine, with the wealth of resources and help from the lab group I was able to eventually overcome any logistic problems or misconceptions with literature. My struggle had another source. There were several times throughout the year where I thought I was approaching the crux of the problem. A certain analysis would show an effect and my mind would autocomplete the follow through on what this means for the larger questions. We are very good at creating stories from sparse evidence. Only for a week later for me to revisit the result, realise I had missed something and the effect I was seeing was a result of a confound or a mistake. A week's worth of cognitive toiling undone and back to square one. But Daksh, surely you would make sure everything is solid, before moving on and using it? Yes you're right but sometimes I genuinely missed things. I try to focus on being more precise, asking "better" questions to my data and not getting ahead of myself. 

There was also many rewarding aspects of being focused deeply on a narrow problem. To test new predictions, I often had to attempt new (or new to me) things methodologically. I have never done any Bayesian statistics, but it was required. So with help of my lab I started at the most basic level, then tried adding different parameters, then tried looking at the model outputs and model fits, then tried writing my own model description and so on. I essentially tried my way through it, learning as I went. Sure, we were trying within a reasonable space, with a broad idea of where things should roughly go but I didn't expect the level of involvement required for analytical decisions. It can seem daunting to work through the space of possible modelling decisions but a lot of the intimidation disappears once you've spent enough time breaking things and fixing them. The relationship with research changes after producing rather than consuming. When I read a paper, my eyes are drawn more to the small decisions and justifications made in methods as well as trying to extract my own interpretations before jumping to the discussion.  

A thesis is never truly finished, it's just at some point it is ready to be shown to the wider world.  As I look towards completing my thesis and my viva, I am focusing on how best to communicate my ideas. Systematically examining all my assumptions, decisions and findings making sure each is backed up with evidence and careful consideration. This thesis feels incredibly personal and I want it to reflect my transition from a consumer of research to a contributor. 


\end{document}
